% This is text associated with the open textbook "Open Electromagnetics"
% See immediately below for author, date
% This work is licensed under the Creative Commons Attribution-ShareAlike 4.0 International License. (CC BY-SA 4.0) 
% To view a copy of the license, visit http://creativecommons.org/licenses/by-sa/4.0/

%ORB TITLE Dielectric Media
%ORB AUTHOR 0        # S.W. Ellingson, Virginia Tech (ellingson@vt.edu)
%ORB DATE 20180806   # yyyymmdd
%ORB ABSTRACT Dielectric Media

\label{m0107_Dielectric_Media} % <-- Must match name of this file and module directory name.  Don't mess with this!
\index{dielectric (material)|(}

\emph{Dielectric} is particular category of materials that exhibit low conductivity\footnote{See Section~\ref{m0007_Properties_of_Materials} for a refresher on parameters describing properties of materials.} because their constituent molecules remain intact when exposed to an electric field, as opposed to shedding electrons as is the case in good conductors.
Subsequently, dielectrics do not effectively pass current, and are therefore considered ``good insulators'' 
\index{insulator!good}
as well as ``poor conductors.'' 
\index{conductor!poor}
An important application of dielectrics in electrical engineering is as a spacer material in printed circuit boards (PCBs), coaxial cables, and capacitors.

Examples of dielectrics include air, 
\index{air}
glass, \index{glass}
teflon, \index{teflon}
 and fiberglass epoxy (the material used in common ``FR4'' 
\index{FR4}
printed circuit boards). 
\index{printed circuit board (PCB)}
These and other dielectrics are listed along with values of their constitutive parameters in Section~\ref{m0135_Table_of_Permittivities}.

The electromagnetic properties of dielectric materials are quantified primarily by relative permittivity 
\index{permittivity!relative}
$\epsilon_r$ (Section~\ref{m0008_Permittivity}), which ranges from very close to 1 upward to roughly 50, and is less than 6 or so for most commonly-encountered materials having low moisture content.  
The permeability 
\index{permeability}
of dielectric materials is approximately equal to the free-space value (i.e., $\mu\approx\mu_0$); therefore, these materials are said to be ``non-magnetic.''
\index{non-magnetic (material)} 

{\bf Additional Reading:}
\begin{itemize}
\item ``\href{https://en.wikipedia.org/wiki/Dielectric}{Dielectric}'' on Wikipedia.
\end{itemize}

\index{dielectric (material)|)}
